% ============================================================
%  高考数学新高考I卷 LaTeX 模板（导言区）
%  综合 2025 卷 (2506.tex) 与 2026 卷 (2606.tex) 排版风格
%  用法：在正文中 % ============================================================
%  高考数学新高考I卷 LaTeX 模板（导言区）
%  综合 2025 卷 (2506.tex) 与 2026 卷 (2606.tex) 排版风格
%  用法：在正文中 % ============================================================
%  高考数学新高考I卷 LaTeX 模板（导言区）
%  综合 2025 卷 (2506.tex) 与 2026 卷 (2606.tex) 排版风格
%  用法：在正文中 % ============================================================
%  高考数学新高考I卷 LaTeX 模板（导言区）
%  综合 2025 卷 (2506.tex) 与 2026 卷 (2606.tex) 排版风格
%  用法：在正文中 \input{gaokao_preamble} 后直接写题目内容
% ============================================================

\documentclass[12pt, a4paper, oneside]{ctexart}

% ── 1. 数学与链接 ──
\usepackage{amsmath, amsthm, amssymb}
\usepackage[bookmarks=true, colorlinks, citecolor=blue, linkcolor=black]{hyperref}

% ── 2. 字体方案（newtxmath Times 风格） ──
\usepackage{newtxmath}

% ── 3. 页面布局（2.5cm 边距） ──
\usepackage[a4paper, margin=2.5cm, footskip=1cm]{geometry}

% ── 4. 页眉页脚 ──
\usepackage{fancyhdr}
\usepackage{lastpage}
\pagestyle{fancy}
\fancyhf{}
\fancyfoot[C]{数学试题第\thepage 页（共\pageref{LastPage}页）}
\renewcommand{\headrulewidth}{0pt}
\renewcommand{\footrulewidth}{0pt}

% ── 5. 表格与插图 ──
\usepackage{graphicx}
\usepackage{adjustbox}
\usepackage{diagbox}
\usepackage{makecell}
\usepackage{caption}
\usepackage{float}
\usepackage{tikz}

% ── 6. 选择题选项（tasks 环境） ──
\usepackage{tasks}
\settasks{
    label       = \Alph*.,
    label-width = 1.8em,
    item-indent = 2.4em,
    label-offset = 0.5em,
    column-sep  = 2em,
    before-skip = 0pt,
    after-skip  = 0pt
}

% ── 7. 大题编号（enumitem 体系） ──
\usepackage[shortlabels]{enumitem}
\newlist{examenum}{enumerate}{3}
\setlist[examenum,1]{label=\arabic*., leftmargin=2em, itemsep=0.3em, parsep=0em}
\setlist[examenum,2]{label=(\arabic*), leftmargin=1.5em, itemsep=0.1em, parsep=0em}
\setlist[examenum,3]{label=(\roman*), leftmargin=1.5em, itemsep=0.1em, parsep=0em}

% ── 8. 自定义命令 ──
\newcommand{\mycircled}[1]{%
  \tikz[baseline=0pt,outer sep=0pt]{%
    \node[draw,circle,inner sep=0.5pt,minimum size=1.4em,
          line width=0.4pt,font=\zihao{-5}] {#1};%
  }%
}
\newcommand{\Parallel}{\raisebox{0.1ex}{\rotatebox[origin=c]{-20}{$\parallel$}}}
\newcommand{\bigstarraised}{\raisebox{0.2ex}{$\bigstar$}}
\newcommand{\blank}{\underline{\hspace{2cm}}}
\newcommand{\subjecttitle}[1]{{\fontsize{24pt}{22pt}\selectfont\centering\textbf{#1}\par}}
\newcommand{\subtitle}[1]{{\fontsize{16pt}{16pt}\selectfont\centering #1\par}}
\renewcommand{\lim}{\limits}
\newcommand{\mi}{\mathrm{i}}
\newcommand{\me}{\mathrm{e}}
 后直接写题目内容
% ============================================================

\documentclass[12pt, a4paper, oneside]{ctexart}

% ── 1. 数学与链接 ──
\usepackage{amsmath, amsthm, amssymb}
\usepackage[bookmarks=true, colorlinks, citecolor=blue, linkcolor=black]{hyperref}

% ── 2. 字体方案（newtxmath Times 风格） ──
\usepackage{newtxmath}

% ── 3. 页面布局（2.5cm 边距） ──
\usepackage[a4paper, margin=2.5cm, footskip=1cm]{geometry}

% ── 4. 页眉页脚 ──
\usepackage{fancyhdr}
\usepackage{lastpage}
\pagestyle{fancy}
\fancyhf{}
\fancyfoot[C]{数学试题第\thepage 页（共\pageref{LastPage}页）}
\renewcommand{\headrulewidth}{0pt}
\renewcommand{\footrulewidth}{0pt}

% ── 5. 表格与插图 ──
\usepackage{graphicx}
\usepackage{adjustbox}
\usepackage{diagbox}
\usepackage{makecell}
\usepackage{caption}
\usepackage{float}
\usepackage{tikz}

% ── 6. 选择题选项（tasks 环境） ──
\usepackage{tasks}
\settasks{
    label       = \Alph*.,
    label-width = 1.8em,
    item-indent = 2.4em,
    label-offset = 0.5em,
    column-sep  = 2em,
    before-skip = 0pt,
    after-skip  = 0pt
}

% ── 7. 大题编号（enumitem 体系） ──
\usepackage[shortlabels]{enumitem}
\newlist{examenum}{enumerate}{3}
\setlist[examenum,1]{label=\arabic*., leftmargin=2em, itemsep=0.3em, parsep=0em}
\setlist[examenum,2]{label=(\arabic*), leftmargin=1.5em, itemsep=0.1em, parsep=0em}
\setlist[examenum,3]{label=(\roman*), leftmargin=1.5em, itemsep=0.1em, parsep=0em}

% ── 8. 自定义命令 ──
\newcommand{\mycircled}[1]{%
  \tikz[baseline=0pt,outer sep=0pt]{%
    \node[draw,circle,inner sep=0.5pt,minimum size=1.4em,
          line width=0.4pt,font=\zihao{-5}] {#1};%
  }%
}
\newcommand{\Parallel}{\raisebox{0.1ex}{\rotatebox[origin=c]{-20}{$\parallel$}}}
\newcommand{\bigstarraised}{\raisebox{0.2ex}{$\bigstar$}}
\newcommand{\blank}{\underline{\hspace{2cm}}}
\newcommand{\subjecttitle}[1]{{\fontsize{24pt}{22pt}\selectfont\centering\textbf{#1}\par}}
\newcommand{\subtitle}[1]{{\fontsize{16pt}{16pt}\selectfont\centering #1\par}}
\renewcommand{\lim}{\limits}
\newcommand{\mi}{\mathrm{i}}
\newcommand{\me}{\mathrm{e}}
 后直接写题目内容
% ============================================================

\documentclass[12pt, a4paper, oneside]{ctexart}

% ── 1. 数学与链接 ──
\usepackage{amsmath, amsthm, amssymb}
\usepackage[bookmarks=true, colorlinks, citecolor=blue, linkcolor=black]{hyperref}

% ── 2. 字体方案（newtxmath Times 风格） ──
\usepackage{newtxmath}

% ── 3. 页面布局（2.5cm 边距） ──
\usepackage[a4paper, margin=2.5cm, footskip=1cm]{geometry}

% ── 4. 页眉页脚 ──
\usepackage{fancyhdr}
\usepackage{lastpage}
\pagestyle{fancy}
\fancyhf{}
\fancyfoot[C]{数学试题第\thepage 页（共\pageref{LastPage}页）}
\renewcommand{\headrulewidth}{0pt}
\renewcommand{\footrulewidth}{0pt}

% ── 5. 表格与插图 ──
\usepackage{graphicx}
\usepackage{adjustbox}
\usepackage{diagbox}
\usepackage{makecell}
\usepackage{caption}
\usepackage{float}
\usepackage{tikz}

% ── 6. 选择题选项（tasks 环境） ──
\usepackage{tasks}
\settasks{
    label       = \Alph*.,
    label-width = 1.8em,
    item-indent = 2.4em,
    label-offset = 0.5em,
    column-sep  = 2em,
    before-skip = 0pt,
    after-skip  = 0pt
}

% ── 7. 大题编号（enumitem 体系） ──
\usepackage[shortlabels]{enumitem}
\newlist{examenum}{enumerate}{3}
\setlist[examenum,1]{label=\arabic*., leftmargin=2em, itemsep=0.3em, parsep=0em}
\setlist[examenum,2]{label=(\arabic*), leftmargin=1.5em, itemsep=0.1em, parsep=0em}
\setlist[examenum,3]{label=(\roman*), leftmargin=1.5em, itemsep=0.1em, parsep=0em}

% ── 8. 自定义命令 ──
\newcommand{\mycircled}[1]{%
  \tikz[baseline=0pt,outer sep=0pt]{%
    \node[draw,circle,inner sep=0.5pt,minimum size=1.4em,
          line width=0.4pt,font=\zihao{-5}] {#1};%
  }%
}
\newcommand{\Parallel}{\raisebox{0.1ex}{\rotatebox[origin=c]{-20}{$\parallel$}}}
\newcommand{\bigstarraised}{\raisebox{0.2ex}{$\bigstar$}}
\newcommand{\blank}{\underline{\hspace{2cm}}}
\newcommand{\subjecttitle}[1]{{\fontsize{24pt}{22pt}\selectfont\centering\textbf{#1}\par}}
\newcommand{\subtitle}[1]{{\fontsize{16pt}{16pt}\selectfont\centering #1\par}}
\renewcommand{\lim}{\limits}
\newcommand{\mi}{\mathrm{i}}
\newcommand{\me}{\mathrm{e}}
 后直接写题目内容
% ============================================================

\documentclass[12pt, a4paper, oneside]{ctexart}

% ── 1. 数学与链接 ──
\usepackage{amsmath, amsthm, amssymb}
\usepackage[bookmarks=true, colorlinks, citecolor=blue, linkcolor=black]{hyperref}

% ── 2. 字体方案（newtxmath Times 风格） ──
\usepackage{newtxmath}

% ── 3. 页面布局（2.5cm 边距） ──
\usepackage[a4paper, margin=2.5cm, footskip=1cm]{geometry}

% ── 4. 页眉页脚 ──
\usepackage{fancyhdr}
\usepackage{lastpage}
\pagestyle{fancy}
\fancyhf{}
\fancyfoot[C]{数学试题第\thepage 页（共\pageref{LastPage}页）}
\renewcommand{\headrulewidth}{0pt}
\renewcommand{\footrulewidth}{0pt}

% ── 5. 表格与插图 ──
\usepackage{graphicx}
\usepackage{adjustbox}
\usepackage{diagbox}
\usepackage{makecell}
\usepackage{caption}
\usepackage{float}
\usepackage{tikz}

% ── 6. 选择题选项（tasks 环境） ──
\usepackage{tasks}
\settasks{
    label       = \Alph*.,
    label-width = 1.8em,
    item-indent = 2.4em,
    label-offset = 0.5em,
    column-sep  = 2em,
    before-skip = 0pt,
    after-skip  = 0pt
}

% ── 7. 大题编号（enumitem 体系） ──
\usepackage[shortlabels]{enumitem}
\newlist{examenum}{enumerate}{3}
\setlist[examenum,1]{label=\arabic*., leftmargin=2em, itemsep=0.3em, parsep=0em}
\setlist[examenum,2]{label=(\arabic*), leftmargin=1.5em, itemsep=0.1em, parsep=0em}
\setlist[examenum,3]{label=(\roman*), leftmargin=1.5em, itemsep=0.1em, parsep=0em}

% ── 8. 自定义命令 ──
\newcommand{\mycircled}[1]{%
  \tikz[baseline=0pt,outer sep=0pt]{%
    \node[draw,circle,inner sep=0.5pt,minimum size=1.4em,
          line width=0.4pt,font=\zihao{-5}] {#1};%
  }%
}
\newcommand{\Parallel}{\raisebox{0.1ex}{\rotatebox[origin=c]{-20}{$\parallel$}}}
\newcommand{\bigstarraised}{\raisebox{0.2ex}{$\bigstar$}}
\newcommand{\blank}{\underline{\hspace{2cm}}}
\newcommand{\subjecttitle}[1]{{\fontsize{24pt}{22pt}\selectfont\centering\textbf{#1}\par}}
\newcommand{\subtitle}[1]{{\fontsize{16pt}{16pt}\selectfont\centering #1\par}}
\renewcommand{\lim}{\limits}
\newcommand{\mi}{\mathrm{i}}
\newcommand{\me}{\mathrm{e}}
